
\let\negmedspace\undefined
\let\negthickspace\undefined
\documentclass[journal]{IEEEtran}
\usepackage[a5paper, margin=10mm, onecolumn]{geometry}
\usepackage{lmodern} % Ensure lmodern is loaded for pdflatex

\setlength{\headheight}{1cm} % Set the height of the header box
\setlength{\headsep}{0mm}     % Set the distance between the header box and the top of the text

\usepackage{gvv-book}
\usepackage{gvv}
\usepackage{cite}
\usepackage{amsmath,amssymb,amsfonts,amsthm}
\usepackage{algorithmic}
\usepackage{graphicx}
\graphicspath{{./figs/}}
\usepackage{textcomp}
\usepackage{xcolor}
\usepackage{txfonts}
\usepackage{listings}
\usepackage{enumitem}
\usepackage{mathtools}
\usepackage{gensymb}
\usepackage{comment}
\usepackage[breaklinks=true]{hyperref}
\usepackage{tkz-euclide} 
\usepackage{listings}
\usepackage{gvv}                                        
\def\inputGnumericTable{}                                 
\usepackage[latin1]{inputenc}                                
\usepackage{color}                                            
\usepackage{array}                                            
\usepackage{longtable}                                       
\usepackage{calc}                                             
\usepackage{multirow}                                         
\usepackage{hhline}                                           
\usepackage{ifthen}                                           
\usepackage{lscape}
\usepackage{circuitikz}
\tikzstyle{block} = [rectangle, draw, fill=blue!20, 
text width=4em, text centered, rounded corners, minimum height=3em]
\tikzstyle{sum} = [draw, fill=blue!10, circle, minimum size=1cm, node distance=1.5cm]
\tikzstyle{input} = [coordinate]
\tikzstyle{output} = [coordinate]



\begin{document}
	
	\bibliographystyle{IEEEtran}
	\vspace{3cm}
	
	\title{1.9.25}
	\author{EE25BTECH11036 - M Chanakya Srinivas}
	\maketitle
	% \newpage
	% \bigskip
	{\let\newpage\relax\maketitle}
	
	\renewcommand{\thefigure}{\theenumi}
	\renewcommand{\thetable}{\theenumi}
	\setlength{\intextsep}{10pt} % Space between text and floats
	
	
\renewcommand\theequation{\arabic{equation}}
\renewcommand\thefigure{\arabic{figure}}
\renewcommand{\thetable}{\theenumi}







\textbf{Question 1.9.25:}  
If the point $P(x,y)$ is equidistant from $A(a+b, b-a)$ and $B(a-b, a+b)$,  
prove that $bx = ay$.

\bigskip
\textbf{Solution:}

\medskip
\textbf{1. Given Data}
\begin{align}
A &= \myvec{a+b \\ b-a}, \quad 
B = \myvec{a-b \\ a+b}, \quad 
P = \myvec{x \\ y} \label{eq:def}
\end{align}

\medskip
\textbf{2. Equidistant Condition}
\begin{align}
\|P-A\|^2 &= \|P-B\|^2 \label{eq:eqdist} \\[6pt]
(P-A)^T(P-A) &= (P-B)^T(P-B) \label{eq:expand1} \\[6pt]
P^TP - 2A^TP + A^TA &= P^TP - 2B^TP + B^TB \label{eq:expand2}
\end{align}

\medskip
\textbf{3. Simplification}
\begin{align}
2(B-A)^TP &= B^TB - A^TA \label{eq:step1} \\[6pt]
B-A &= \myvec{-2b \\ 2a} = 2\myvec{-b \\ a} \label{eq:diff} \\[6pt]
A^TA &= (a+b)^2 + (b-a)^2 = 2a^2 + 2b^2 \label{eq:ATA} \\[6pt]
B^TB &= (a-b)^2 + (a+b)^2 = 2a^2 + 2b^2 \label{eq:BTB} \\[6pt]
B^TB - A^TA &= 0 \label{eq:rhs0} \\[6pt]
(B-A)^TP &= 0 \label{eq:lhs} \\[6pt]
\myvec{-2b & 2a}\myvec{x \\ y} &= 0 \label{eq:prod}
\end{align}

\medskip
\textbf{4. Final Result}
\begin{align}
-2bx + 2ay &= 0 
\quad \Rightarrow \quad 
\boxed{bx = ay} \label{eq:final}
\end{align}

\noindent
Thus, the locus of $P$ is the line
\begin{align*}
    n^Tx = 0, \quad n = \myvec{-b \\ a}.
\end{align*}
\begin{figure}
       \centering
    \includegraphics[width=0.6\columnwidth]{figs/Figure1.png}
    \caption{}
    \label{fig:placeholder}
\end{figure}
\begin{figure}
    \centering
    \includegraphics[width=0.8\columnwidth]{figs/Figure_4.png}
    \caption{}
    \label{fig:placeholder}
\end{figure}
\end{document}